\documentclass{situationist-preprint}

\title{Doxastic Involuntarism and Attention}
\msid{EPI-11-a}
\msfamily{EPI-11}
\msdiscipline{Epistemology}
\mstype{Principal manuscript}
\mskeywords{doxastic voluntarism; belief; attention; acceptance; epistemic responsibility; deontology of belief}

\begin{document}
\maketitle

\begin{abstract}
Can one believe at will? Descartes held that judgment is an act of the will, and some contemporary epistemologists hold that we have a voluntary control over belief that grounds our responsibility for it. We argue that belief is not under direct voluntary control, but that attention is, and that belief follows attention. The first claim rests on the nature of belief: a state counts as a belief that $p$ only if it is formed in response to considerations taken to bear on whether $p$, while a decision responds to considerations bearing on whether to do something. A state produced directly by a decision is therefore not a belief. The second claim is a datum about attention: we can decide what to look at, what question to consider and what evidence to review. The third is a principle of rational functioning: the attitude a rational agent takes toward $p$ is fixed by the evidence she attends to. Together they give us an indirect control over belief, bounded by the available evidence, which is enough to ground the deontological vocabulary of epistemic evaluation and locates our responsibility for belief in our responsibility for inquiry.
\end{abstract}

\section{The problem}

Descartes, in the Fourth Meditation, explained error by the will's extending further than the intellect: judgment is an act of the will, and we err when we affirm what we do not clearly perceive. If judgment is voluntary, so, it seems, is belief, and the ethics of belief is an ethics of a kind of action. \citet{Williams1973} argued against this that believing at will is impossible, and not merely difficult, because a state one knew to have been produced by one's will regardless of the truth could not be recognized by one as a belief. \citet{Bennett1990} and \citet{Alston1988} defended the weaker claim that we in fact have no such power, and Alston drew the consequence that the deontological conception of epistemic justification, on which justified belief is belief in which one has done one's epistemic duty, cannot be sustained, since \enquote{ought} implies \enquote{can}.

Voluntarists have replied in two ways. \citet{Ginet2001} gave cases in which we seem to decide to believe, as when one decides to believe that one locked the door in order to stop worrying about it. \citet{Steup2008} argued that belief is free in the sense that matters for responsibility, since it is responsive to reasons, and that the involuntarist demands of belief a kind of control that is not required of action either. On the other side, \citet{Feldman2001} argued that epistemic evaluation does not need voluntary control at all: it describes what fits the evidence, not what one is obliged to do.

We argue that belief is not under direct voluntary control, that attention is, and that belief follows attention. The result vindicates Alston's premise and rejects his conclusion: there is a deontology of belief, but its objects are acts of attention and inquiry.

\section{Framework}

\begin{definition}[Direct and indirect control]\label{d:control}
An agent has \emph{direct voluntary control} over whether she is in a state $S$ if she can bring it about that she is in $S$ by deciding to, without doing something else as a means. She has \emph{indirect voluntary control} over $S$ if she can bring it about by doing something else over which she has direct voluntary control.
\end{definition}

\begin{principle}[Belief answers to truth]\label{p:truth}
A state is a belief that $p$ only if it is formed and maintained in response to considerations that the subject takes to bear on whether $p$ is true.
\end{principle}

Principle~\ref{p:truth} is a version of the thesis that belief constitutively aims at truth, in the form defended by \citet{ShahVelleman2005}: when one deliberates whether to believe $p$, the question gives way to the question whether $p$, and only considerations bearing on the latter can settle it. It is also what makes self-knowledge of belief possible by the transparency method of EPI-15: one knows whether one believes $p$ by considering whether $p$, which would be no guide if beliefs could be formed in response to other considerations.

\begin{assumption}[Decisions answer to value]\label{a:decide}
A decision to bring about a state is made in response to considerations that the agent takes to bear on whether to bring it about, that is, on its value or the value of its effects.
\end{assumption}

\begin{assumption}[Attention is voluntary]\label{a:attention}
An agent has direct voluntary control over what she attends to: which question she considers, where she looks, which evidence she reviews and for how long.
\end{assumption}

Assumption~\ref{a:attention} records a familiar fact, central to the psychology of attention since \citet[ch.~XI]{James1890}, who held that the effort of attention is the essential phenomenon of the will. Not all attention is voluntary, since attention is also captured by salient stimuli, but voluntary attention is the paradigm of a mental act that one performs by deciding to.

\begin{assumption}[Belief follows attention]\label{a:follows}
For a rational agent, the doxastic attitude toward $p$ is fixed by the evidence bearing on $p$ to which she attends, together with her background beliefs.
\end{assumption}

The attitude of Assumption~\ref{a:follows} need not be the one that fits the agent's total evidence. By EPI-21, one's evidence is one's knowledge, and one knows much that one is not attending to. An agent whose attention is partial forms the attitude that the attended part of her evidence supports, and this may differ from the attitude that her total evidence supports, which by EPI-17 is unique.

\section{The argument}

\begin{theorem}[Involuntarism]\label{t:invol}
Belief is not under direct voluntary control.
\end{theorem}

\begin{proof}
Suppose an agent brings it about that she is in a state $S$ by deciding to, without doing anything else as a means. By Assumption~\ref{a:decide}, her decision responds to considerations bearing on whether to be in $S$. Since nothing intervenes between the decision and the state, $S$ is formed in response to those considerations and to nothing else. If $S$ were a belief that $p$, then by Principle~\ref{p:truth} it would be formed in response to considerations that she takes to bear on whether $p$ is true. Considerations about whether to believe $p$ bear on whether $p$ is true only in exceptional cases, such as those in which believing $p$ makes $p$ true, and in those cases what forms the state is the agent's recognition of the evidence, not her decision. So $S$ is not a belief that $p$.
\end{proof}

The proof explains the phenomenon that Pascal's interlocutor reports, on which REL-06-a relies in reading the wager as a recommendation of commitment: offered the best of practical reasons for believing, he finds that he cannot. Practical reasons are reasons for deciding, and deciding is the wrong kind of response to produce a belief. This is the conceptual point of \citet{Williams1973} without its reliance on what the subject knows about the origin of her state: a state produced directly by a decision is not a belief, whether or not the subject is aware of how it was produced.

\begin{theorem}[Indirect control]\label{t:indirect}
A rational agent has indirect voluntary control over her beliefs, through attention. The control is bounded: she can bring herself to believe $p$ only if some body of available evidence to which she could attend supports $p$.
\end{theorem}

\begin{proof}
By Assumption~\ref{a:attention}, she can decide which evidence to attend to. By Assumption~\ref{a:follows}, the attitude she then takes toward $p$ is fixed by that evidence and her background. Hence, by deciding what to attend to, she brings about her attitude toward $p$, which is indirect control in the sense of Definition~\ref{d:control}. If no body of evidence available for her attention supports $p$, no decision about attention leads to belief that $p$, by Assumption~\ref{a:follows}.
\end{proof}

\begin{corollary}[Responsibility]\label{c:resp}
Epistemic obligations that require an agent to believe or not to believe are satisfiable, and they are satisfied or violated through what she does with her attention: through what she considers, investigates and ignores.
\end{corollary}

Corollary~\ref{c:resp} answers Alston's argument. Alston was right that there can be no obligation to perform an act of believing, since there is no such act. But an obligation can bear on a state that one controls indirectly, as one may be obliged to be on time, which one achieves by leaving early. The deontology of belief is a deontology of inquiry. It is the form in which REL-17-a and REL-06-a use the present result: practical reasons can recommend acts and commitments, including commitments that direct attention, but cannot directly recommend beliefs; and a belief that results from such attention is assessed by the evidence attended to, not by the reasons for attending.

\section{Objections}

\begin{objection}[Deciding to believe]
I have checked the door once and cannot remember whether I locked it. I decide to believe that I did, and go on my way. I have decided to believe \citep{Ginet2001}.
\end{objection}

\begin{reply}
What one decides in such a case is to stop inquiring and to act on the proposition, taking it as settled for the purposes at hand. This is acceptance in the sense of \citet{Cohen1992} and \citet{Bratman1992}: a policy of using a proposition as a premise in practical reasoning, which is voluntary and need not be accompanied by belief. One may accept that the door is locked while still having the nagging doubt, which shows that the decision did not produce belief. Where belief does follow, it follows the evidence one has, in this case one's memory of a routine, once one stops attending to the possibility of an exception, as Theorem~\ref{t:indirect} predicts.
\end{reply}

\begin{objection}[Doxastic compatibilism]
Belief is free in the sense that matters: it is responsive to reasons, and the agent can be held responsible for it. Voluntary control of the kind defined is not required even for action; the involuntarist sets an impossible standard \citep{Steup2008}.
\end{objection}

\begin{reply}
We grant that belief is responsive to reasons, and that this is a form of agency, which \citet{Hieronymi2006} calls evaluative control: in settling whether $p$ one settles what one believes. But it is not voluntary control, since what settles the matter is the evidence and not a decision, and action differs in just this respect: one can perform an action for any reason that one takes to favour performing it. Theorem~\ref{t:invol} denies only direct voluntary control; it is compatible with evaluative control, and Corollary~\ref{c:resp} supplies the voluntary control that the deontology of belief needs.
\end{reply}

\begin{objection}[Belief does not follow attention]
People attend to strong evidence and continue to believe what it refutes, and phobic or self-deceived subjects keep their attitudes whatever they attend to. Assumption~\ref{a:follows} is false.
\end{objection}

\begin{reply}
Assumption~\ref{a:follows} concerns rational agents, and the cases described are failures of rationality. They are failures of exactly the kind that EPI-15 identifies as the limits of transparency, attitudes that are not answerable to reasons; and to the extent that an attitude is not answerable to reasons, it is not controlled by attention either, so that Theorem~\ref{t:indirect} is restricted in the same way. This is the correct result: the recalcitrant phobic cannot think her way out of her fear, and so cannot attend her way out of it either.
\end{reply}

\begin{objection}[Voluntarism by transitivity]
If attention is voluntary and belief follows attention, then belief is voluntary after all, and the agent who selectively attends to the evidence for a comforting conclusion believes it at will.
\end{objection}

\begin{reply}
She believes it by means of selective attention, which is indirect control, and the difference matters. The control is bounded by the evidence available (Theorem~\ref{t:indirect}), so that no amount of attention produces belief in what nothing supports. It requires knowing in advance which evidence favours the conclusion. And it leaves the belief open to assessment by the total evidence, which by EPI-21 includes what she knows and chose not to consider. Selective attention of this kind is a vice, the one that Clifford's shipowner exhibits when he puts his doubts aside, and that EPI-23 examines; Corollary~\ref{c:resp} locates the wrong in the attention, which is where it lies.
\end{reply}

\section{Conclusion}

A state produced directly by a decision responds to reasons for being in it, while a belief responds to reasons for thinking it true; so belief is not under direct voluntary control. Attention is, and a rational agent's beliefs are fixed by the evidence she attends to. We therefore control our beliefs indirectly, within the limits of the evidence available to us, and our epistemic obligations are obligations about what we consider, investigate and ignore.

\begin{thebibliography}{99}\small
\bibitem[Alston(1988)]{Alston1988} Alston, W.~P. (1988). The deontological conception of epistemic justification. \emph{Philosophical Perspectives}, 2, 257--299.
\bibitem[Bennett(1990)]{Bennett1990} Bennett, J. (1990). Why is belief involuntary? \emph{Analysis}, 50(2), 87--107.
\bibitem[Bratman(1992)]{Bratman1992} Bratman, M.~E. (1992). Practical reasoning and acceptance in a context. \emph{Mind}, 101(401), 1--15.
\bibitem[Cohen(1992)]{Cohen1992} Cohen, L.~J. (1992). \emph{An Essay on Belief and Acceptance}. Clarendon Press.
\bibitem[Feldman(2001)]{Feldman2001} Feldman, R. (2001). Voluntary belief and epistemic evaluation. In M. Steup (ed.), \emph{Knowledge, Truth, and Duty: Essays on Epistemic Justification, Responsibility, and Virtue}. Oxford University Press.
\bibitem[Ginet(2001)]{Ginet2001} Ginet, C. (2001). Deciding to believe. In M. Steup (ed.), \emph{Knowledge, Truth, and Duty: Essays on Epistemic Justification, Responsibility, and Virtue}. Oxford University Press.
\bibitem[Hieronymi(2006)]{Hieronymi2006} Hieronymi, P. (2006). Controlling attitudes. \emph{Pacific Philosophical Quarterly}, 87(1), 45--74.
\bibitem[James(1890)]{James1890} James, W. (1890). \emph{The Principles of Psychology}. Henry Holt.
\bibitem[Shah and Velleman(2005)]{ShahVelleman2005} Shah, N. and Velleman, J.~D. (2005). Doxastic deliberation. \emph{The Philosophical Review}, 114(4), 497--534.
\bibitem[Steup(2008)]{Steup2008} Steup, M. (2008). Doxastic freedom. \emph{Synthese}, 161(3), 375--392.
\bibitem[Williams(1973)]{Williams1973} Williams, B. (1973). Deciding to believe. In \emph{Problems of the Self}. Cambridge University Press.
\end{thebibliography}

\end{document}
